\documentclass[12pt,a4paper]{article}
\usepackage[utf8]{inputenc}
\usepackage[T1]{fontenc}
\usepackage{amsmath,amssymb,amsthm}
\usepackage{graphicx}
\usepackage{listings}
\usepackage{xcolor}
\usepackage{hyperref}
\usepackage{geometry}
\usepackage{tcolorbox}
\usepackage{needspace}
\usepackage{tikz}
\usepackage{array}
\usepackage{booktabs}
\usepackage{pifont}
\usetikzlibrary{arrows,positioning,matrix}
\geometry{margin=2.5cm}

% Command for circled numbers
\newcommand*\circled[1]{\tikz[baseline=(char.base)]{\node[shape=circle,draw,inner sep=2pt] (char) {#1};}}

% Couleurs personnalisees
\definecolor{titrebleu}{RGB}{0,102,204}
\definecolor{defvert}{RGB}{0,153,76}
\definecolor{exempleorange}{RGB}{255,102,0}
\definecolor{algoviolet}{RGB}{153,0,153}
\definecolor{importantrouge}{RGB}{204,0,0}
\definecolor{fondgris}{RGB}{245,245,245}
\definecolor{fondjaune}{RGB}{255,250,205}
\definecolor{fondvert}{RGB}{230,255,230}
\definecolor{fondbleu}{RGB}{230,240,255}
\definecolor{fondrose}{RGB}{255,230,240}

% Boites colorees
\newtcolorbox{definition}{
    colback=fondvert,
    colframe=defvert,
    fonttitle=\bfseries,
    title=Definition,
    arc=3mm
}

\newtcolorbox{exemple}{
    colback=fondjaune,
    colframe=exempleorange,
    fonttitle=\bfseries,
    title=Exemple,
    arc=3mm
}

\newtcolorbox{important}{
    colback=white,
    colframe=importantrouge,
    fonttitle=\bfseries,
    title=Important,
    arc=3mm
}

\newtcolorbox{astuce}{
    colback=fondbleu,
    colframe=titrebleu,
    fonttitle=\bfseries,
    title=A savoir,
    arc=3mm
}

\newtcolorbox{methode}{
    colback=fondbleu,
    colframe=algoviolet,
    fonttitle=\bfseries,
    title=Methode,
    arc=3mm
}

\newtcolorbox{notecours}{
    colback=fondrose,
    colframe=importantrouge,
    fonttitle=\bfseries,
    title=Note de Cours,
    arc=3mm
}

\newtcolorbox{theoreme}{
    colback=fondgris,
    colframe=algoviolet,
    fonttitle=\bfseries,
    title=Theoreme,
    arc=3mm
}

% Style des sections
\usepackage{titlesec}
\titleformat{\section}
{\color{titrebleu}\normalfont\Large\bfseries}
{\color{titrebleu}\thesection}{1em}{}

\titleformat{\subsection}
{\color{defvert}\normalfont\large\bfseries}
{\color{defvert}\thesubsection}{1em}{}

\title{\textbf{\Huge\color{titrebleu}La Theorie de la Dualite}\\[0.5em]
\large\color{defvert}Probleme Dual, Theoremes de Dualite et Ecarts Complementaires\\
\small Modeles Lineaires de la Recherche Operationnelle}
\author{\large Prof. Herve KERIVIN, PhD\\
\small Universite Clermont Auvergne\\
Clermont Auvergne INP - ISIMA\\
LIMOS UMR 6158 CNRS}
\date{\large Version 2024-2025 - Cours 4}


\begin{document}

\maketitle
\tableofcontents
\newpage

%=============================================
\section{Motivation : Bornes Superieures}
%=============================================

\needspace{10\baselineskip}
\subsection{Probleme de Depart}

\begin{notecours}
\textbf{Rappel -- La methode du simplexe produit une solution optimale, mais peut-on la verifier sans la relancer ?}

La methode du simplexe nous donne une solution realisable qui maximise $z$. Mais comment \textcolor{importantrouge}{\textbf{prouver}} qu'une solution donnee est optimale sans executer tout l'algorithme ?
\end{notecours}

\begin{exemple}
\textbf{Probleme de PL de reference :}

\begin{align*}
\text{maximiser} \quad & z = 4x_1 + x_2 + 5x_3 + 3x_4\\
\text{sujet a} \quad & x_1 - x_2 - x_3 + 3x_4 \leq 1\\
& 5x_1 + x_2 + 3x_3 + 8x_4 \leq 55\\
& -x_1 + 2x_2 + 3x_3 - 5x_4 \leq 3\\
& x_1, x_2, x_3, x_4 \geq 0
\end{align*}

\textbf{Observations initiales :}
\begin{itemize}
    \item Toute solution realisable donne une \textcolor{defvert}{\textbf{borne inferieure}} sur $z^*$
    \item Par exemple : $x_1 = 3, x_2 = 0, x_3 = 2, x_4 = 0$ donne $z \geq 22$
    \item Encore mieux : $x_1 = 0, x_2 = 14, x_3 = 0, x_4 = 5$ donne $z \geq 29$
\end{itemize}

\textbf{Probleme :} Pour prouver qu'une solution est optimale, il faut aussi une \textcolor{importantrouge}{\textbf{borne superieure}} sur $z^*$.
\end{exemple}

\subsection{Construction d'une Borne Superieure}

\begin{methode}
\textbf{Idee fondamentale -- Combinaison lineaire des contraintes :}

On peut multiplier une contrainte par un scalaire positif pour obtenir une borne.

\textbf{Exemple 1 :} Multiplier la 2e contrainte par $\frac{5}{3}$ :

$$\frac{5}{3}(5x_1 + x_2 + 3x_3 + 8x_4 \leq 55) \quad \Rightarrow \quad \frac{25}{3}x_1 + \frac{5}{3}x_2 + 5x_3 + \frac{40}{3}x_4 \leq \frac{275}{3}$$

Puisque pour toute solution realisable :
$$4x_1 \leq \frac{25}{3}x_1, \quad x_2 \leq \frac{5}{3}x_2, \quad 5x_3 \leq 5x_3, \quad 3x_4 \leq \frac{40}{3}x_4$$

on obtient :
$$z = 4x_1 + x_2 + 5x_3 + 3x_4 \leq \frac{25}{3}x_1 + \frac{5}{3}x_2 + 5x_3 + \frac{40}{3}x_4 \leq \frac{275}{3}$$

Donc $z^* \leq \frac{275}{3} \approx 91{,}7$.

\textbf{Exemple 2 :} Additionner les 2e et 3e contraintes :
$$4x_1 + 3x_2 + 6x_3 + 3x_4 \leq 58 \quad \Rightarrow \quad z^* \leq 58$$

\textbf{Peut-on trouver une borne encore plus fine ?}
\end{methode}

\newpage
\subsection{Combinaison Lineaire Generale}

\begin{methode}
\textbf{Construction systematique :}

On forme une combinaison lineaire des $m$ contraintes avec des multiplicateurs $y_1, y_2, y_3 \geq 0$ :

$$y_1(x_1 - x_2 - x_3 + 3x_4 \leq 1), \quad y_2(5x_1 + x_2 + 3x_3 + 8x_4 \leq 55), \quad y_3(-x_1 + 2x_2 + 3x_3 - 5x_4 \leq 3)$$

Toute solution realisable satisfait alors :
$$(y_1 + 5y_2 - y_3)x_1 + (-y_1 + y_2 + 2y_3)x_2 + (-y_1 + 3y_2 + 3y_3)x_3 + (3y_1 + 8y_2 - 5y_3)x_4 \leq y_1 + 55y_2 + 3y_3$$

Pour conclure que $z^* \leq y_1 + 55y_2 + 3y_3$, il faut que chaque coefficient de la combinaison soit $\geq$ au coefficient de $z$ :
\begin{align*}
y_1 + 5y_2 - y_3 &\geq 4\\
-y_1 + y_2 + 2y_3 &\geq 1\\
-y_1 + 3y_2 + 3y_3 &\geq 5\\
3y_1 + 8y_2 - 5y_3 &\geq 3
\end{align*}
\end{methode}

\begin{important}
\textbf{Conclusion fondamentale :}

Pour obtenir la \textcolor{importantrouge}{\textbf{borne superieure la plus fine possible}}, il faut \textcolor{defvert}{\textbf{minimiser}} $y_1 + 55y_2 + 3y_3$ sous les contraintes ci-dessus avec $y_1, y_2, y_3 \geq 0$.

\textcolor{titrebleu}{\textbf{Ce nouveau probleme d'optimisation est appele le probleme dual !}}
\end{important}

%=============================================
\section{Le Probleme Dual}
%=============================================

\needspace{10\baselineskip}
\subsection{Definition Formelle}

\begin{definition}
\textbf{Probleme primal et probleme dual :}

Considerons le \textcolor{titrebleu}{\textbf{probleme primal}} sous forme standard :
\begin{align*}
\text{maximiser} \quad & z = \sum_{j=1}^{n} c_j x_j\\
\text{sujet a} \quad & \sum_{j=1}^{n} a_{ij} x_j \leq b_i \quad \text{pour } i = 1, 2, \ldots, m\\
& x_j \geq 0 \quad \text{pour } j = 1, 2, \ldots, n
\end{align*}

Son \textcolor{defvert}{\textbf{probleme dual}} est :
\begin{align*}
\text{minimiser} \quad & w = \sum_{i=1}^{m} b_i y_i\\
\text{sujet a} \quad & \sum_{i=1}^{m} a_{ij} y_i \geq c_j \quad \text{pour } j = 1, 2, \ldots, n\\
& y_i \geq 0 \quad \text{pour } i = 1, 2, \ldots, m
\end{align*}
\end{definition}

\begin{astuce}
\textbf{Correspondances primal/dual -- A retenir :}

\begin{center}
\begin{tabular}{l|l}
\toprule
\textbf{PRIMAL} & \textbf{DUAL} \\
\midrule
maximisation & minimisation \\
contrainte $\leq$ & variable $\geq 0$ \\
contrainte $\geq$ & variable $\leq 0$ \\
contrainte $=$ & variable libre \\
variable $\geq 0$ & contrainte $\geq$ \\
variable $\leq 0$ & contrainte $\leq$ \\
variable libre & contrainte $=$ \\
membre droit $b_i$ & coefficient dans l'objectif dual \\
coefficient dans l'objectif $c_j$ & membre droit dual \\
\bottomrule
\end{tabular}
\end{center}
\end{astuce}

\newpage
\subsection{Correspondance Contraintes-Variables}

\begin{notecours}
\textbf{Comment lire la correspondance primal/dual :}

\begin{itemize}
    \item Les $m$ \textcolor{titrebleu}{\textbf{contraintes primales}} $\sum_{j=1}^n a_{ij}x_j \leq b_i$ sont en correspondance \textbf{biunivoque} avec les $m$ \textcolor{defvert}{\textbf{variables duales}} $y_i$
    \item Les $n$ \textcolor{titrebleu}{\textbf{variables primales}} $x_j$ sont en correspondance \textbf{biunivoque} avec les $n$ \textcolor{defvert}{\textbf{contraintes duales}} $\sum_{i=1}^m a_{ij}y_i \geq c_j$
    \item Le coefficient $c_j$ de l'objectif primal devient le \textcolor{importantrouge}{\textbf{membre droit}} de la contrainte duale associee a $x_j$
    \item Le membre droit $b_i$ de la contrainte primale devient le \textcolor{importantrouge}{\textbf{coefficient}} de $y_i$ dans l'objectif dual
\end{itemize}
\end{notecours}

\begin{theoreme}
\textbf{Proposition 1 -- Auto-dualite :}

\textit{Le dual du dual est le primal.}
\end{theoreme}

\begin{definition}
\textbf{Solutions primale et duale :}

\begin{itemize}
    \item Une solution realisable du probleme primal est appelee une \textcolor{titrebleu}{\textbf{solution primale}}
    \item Une solution realisable du probleme dual est appelee une \textcolor{defvert}{\textbf{solution duale}}
\end{itemize}
\end{definition}

\subsection{Exemples de Construction du Dual}

\begin{exemple}
\textbf{Probleme \#1 -- Primal avec 7 contraintes $\leq$ et maximisation :}

\begin{align*}
\text{maximiser} \quad & z = 4x_1 + 3x_2 + 7x_3 + 9x_4\\
\text{sujet a} \quad & 2x_1 - 4x_2 + x_3 \leq -1\\
& x_1 + 5x_2 + x_3 + x_4 \leq 16\\
& x_1 + x_3 \leq 5\\
& 2x_1 + 4x_2 - x_4 \leq 8\\
& x_2 - 3x_3 + x_4 \leq 0\\
& -4x_2 + 3x_4 \leq 4\\
& 5x_1 + 2x_2 - 3x_3 + 6x_4 \leq 19\\
& x_1, x_2, x_3, x_4 \geq 0
\end{align*}

\textbf{Dual correspondant :}

Les variables primales sont non negatives et le probleme est une maximisation : les contraintes duales sont donc des contraintes $\geq$.

\begin{align*}
\text{minimiser} \quad & w = -y_1 + 16y_2 + 5y_3 + 8y_4 + 0 \cdot y_5 + 4y_6 + 19y_7\\
\text{sujet a} \quad & 2y_1 + y_2 + y_3 + 2y_4 + 5y_7 \geq 4\\
& -4y_1 + 5y_2 + 4y_4 + y_5 - 4y_6 + 2y_7 \geq 3\\
& y_1 + y_2 + y_3 - 3y_5 - 3y_7 \geq 7\\
& y_2 - y_4 + y_5 + 3y_6 + 6y_7 \geq 9\\
& y_i \geq 0 \quad \text{pour } i = 1, \ldots, 7
\end{align*}
\end{exemple}

\newpage

\begin{exemple}
\textbf{Probleme \#2 -- Primal avec contraintes $\geq$ et minimisation :}

\begin{align*}
\text{minimiser} \quad & z = 5x_1 + 2x_2 + 6x_3\\
\text{sujet a} \quad & 2x_1 + x_2 + x_3 \geq 5\\
& x_2 + 2x_3 \geq 4\\
& x_1 + x_3 \geq 4\\
& x_1, x_2, x_3 \geq 0
\end{align*}

Les variables primales sont non negatives et le probleme est une minimisation : les contraintes duales sont des contraintes $\leq$.

\begin{align*}
\text{maximiser} \quad & w = 5y_1 + 4y_2 + 4y_3\\
\text{sujet a} \quad & 2y_1 + y_3 \leq 5\\
& y_1 + y_2 \leq 2\\
& y_1 + 2y_2 + y_3 \leq 6\\
& y_i \geq 0 \quad \text{pour } i = 1, 2, 3
\end{align*}
\end{exemple}

\begin{exemple}
\textbf{Probleme \#3 -- Primal avec contraintes mixtes et variables libres :}

\begin{align*}
\text{maximiser} \quad & z = 3x_1 + 2x_2 + 5x_3\\
\text{sujet a} \quad & 5x_1 + 3x_2 + x_3 = -8\\
& 4x_1 + 2x_2 + 8x_3 \leq 1\\
& 6x_1 + 7x_2 + 3x_3 \geq 1\\
& x_1 \leq 4\\
& x_3 \geq 0 \quad \text{(seule contrainte de signe)}
\end{align*}

\textbf{Correspondance :}
\begin{itemize}
    \item Contrainte $=$ $\Rightarrow$ variable duale $y_1$ \textcolor{algoviolet}{\textbf{libre}}
    \item Contrainte $\leq$ $\Rightarrow$ variable duale $y_2 \geq 0$
    \item Contrainte $\geq$ $\Rightarrow$ variable duale $y_3 \leq 0$
    \item $x_1, x_2$ libres $\Rightarrow$ contraintes duales correspondantes sont des \textcolor{algoviolet}{\textbf{egalites}}
    \item $x_3 \geq 0$ $\Rightarrow$ contrainte duale correspondante est une \textcolor{defvert}{\textbf{inegalite $\geq$}}
\end{itemize}
\end{exemple}

%=============================================
\section{Theoremes de Dualite}
%=============================================

\needspace{10\baselineskip}
\subsection{Theoreme de Dualite Faible}

\begin{theoreme}
\textbf{Theoreme 2 (Dualite faible) :}

\textit{Supposons que $x$ soit une solution primale et $y$ une solution duale. Alors :}
$$\sum_{j=1}^{n} c_j x_j \leq \sum_{i=1}^{m} b_i y_i$$
\end{theoreme}

\begin{methode}
\textbf{Preuve du theoreme de dualite faible :}

Puisque $y$ est une solution duale, $\sum_{i=1}^m a_{ij} y_i \geq c_j$ pour tout $j$. Comme $x_j \geq 0$ :
$$\sum_{j=1}^n c_j x_j \leq \sum_{j=1}^n \left(\sum_{i=1}^m a_{ij} y_i\right) x_j = \sum_{i=1}^m \left(\sum_{j=1}^n a_{ij} x_j\right) y_i \leq \sum_{i=1}^m b_i y_i$$

La derniere inegalite vient du fait que $x$ est une solution primale ($\sum_j a_{ij} x_j \leq b_i$) et $y_i \geq 0$. $\square$
\end{methode}

\begin{astuce}
\textbf{Interpretation geometrique :}

Le theoreme de dualite faible dit que \textcolor{importantrouge}{\textbf{toute solution duale fournit une borne superieure}} sur la valeur de la fonction objectif primale, et \textcolor{defvert}{\textbf{toute solution primale fournit une borne inferieure}} sur la valeur de la fonction objectif duale.

Autrement dit, les valeurs de $z$ primal et $w$ dual sont \textcolor{titrebleu}{\textbf{separees}} : valeurs primales $\leq$ valeurs duales.
\end{astuce}

\newpage
\subsection{Corollaire : Critere d'Optimalite par la Dualite}

\begin{theoreme}
\textbf{Corollaire 3 :}

\textit{Supposons que $x^*$ soit une solution primale, $y^*$ une solution duale et :}
$$\sum_{j=1}^{n} c_j x_j^* = \sum_{i=1}^{m} b_i y_i^*$$
\textit{Alors $x^*$ est une solution primale optimale, et $y^*$ est une solution duale optimale.}
\end{theoreme}

\begin{exemple}
\textbf{Application au probleme de reference :}

On peut verifier que les solutions suivantes sont optimales :

\textbf{Solution primale :} $x_1^* = 0, x_2^* = 14, x_3^* = 0, x_4^* = 5$

$$z^* = 4(0) + 1(14) + 5(0) + 3(5) = 29$$

\textbf{Solution duale :} $y_1^* = 11, y_2^* = 0, y_3^* = 6$

$$w^* = 1(11) + 55(0) + 3(6) = 29$$

Comme $z^* = w^* = 29$, les deux solutions sont \textcolor{defvert}{\textbf{simultanement optimales}} !
\end{exemple}

\subsection{Comment Lire la Solution Duale dans le Dictionnaire Primal}

\begin{methode}
\textbf{Lecture de la solution duale optimale a partir du dictionnaire primal optimal :}

\textbf{Dictionnaire optimal pour le probleme primal :}
\begin{align*}
x_2 &= 14 - 2x_1 - 4x_3 - 5x_5 - 3x_7\\
x_4 &= 5 - x_1 - x_3 - 2x_5 - x_7\\
x_6 &= 1 + 5x_1 + 9x_3 + 21x_5 + 11x_7\\
z &= 29 - x_1 - 2x_3 - 11x_5 - 6x_7
\end{align*}

ou $x_5, x_6, x_7$ sont les variables d'ecart des trois contraintes primales.

\textbf{Regle de lecture :}

Les variables d'ecart $x_5, x_6, x_7$ correspondent aux variables duales $y_1, y_2, y_3$ :
$$x_5 \leftrightarrow y_1, \quad x_6 \leftrightarrow y_2, \quad x_7 \leftrightarrow y_3$$

La solution duale optimale s'obtient en \textcolor{importantrouge}{\textbf{prenant les coefficients opposes des variables d'ecart dans l'equation de $z$}} :
$$y_1^* = 11, \quad y_2^* = 0, \quad y_3^* = 6$$
\end{methode}

\begin{important}
\textbf{Pourquoi cette correspondance fonctionne-t-elle ?}

Dans le dictionnaire optimal, le coefficient de la variable d'ecart $x_{n+i}$ dans $z$ est $-\bar{c}_{n+i}$.

La dualite forte (voir section suivante) garantit que ces valeurs correspondent exactement a la solution duale optimale.

\textcolor{titrebleu}{\textbf{La methode du simplexe resout simultanement le primal et le dual !}}
\end{important}

\newpage
\subsection{Theoreme de Dualite Forte}

\begin{theoreme}
\textbf{Theoreme 4 (Dualite forte) :}

\textit{Si le probleme primal a une solution optimale $(x_1^*, x_2^*, \ldots, x_n^*)$, alors le probleme dual a une solution optimale $(y_1^*, y_2^*, \ldots, y_m^*)$ telle que :}
$$\sum_{j=1}^{n} c_j x_j^* = \sum_{i=1}^{m} b_i y_i^*$$
\end{theoreme}

\begin{theoreme}
\textbf{Corollaire 5 :}

\begin{enumerate}
    \item[\circled{i}] \textit{Le probleme primal a une solution optimale \textcolor{defvert}{\textbf{si et seulement si}} le probleme dual a une solution optimale.}
    \item[\circled{ii}] \textit{Si le primal est \textcolor{importantrouge}{\textbf{non borne}}, alors le dual est \textcolor{importantrouge}{\textbf{infaisable}}.}
    \item[\circled{iii}] \textit{Si le dual est \textcolor{importantrouge}{\textbf{non borne}}, alors le primal est \textcolor{importantrouge}{\textbf{infaisable}}.}
\end{enumerate}
\end{theoreme}

\subsection{Tableau des Combinaisons Primal/Dual}

\begin{important}
\textbf{Toutes les combinaisons possibles entre primal et dual :}

\begin{center}
\begin{tabular}{l|c|c|c}
\toprule
& \textbf{Dual Optimal} & \textbf{Dual Infaisable} & \textbf{Dual Non Borne} \\
\midrule
\textbf{Primal Optimal} & \textcolor{defvert}{Possible} & Impossible & Impossible \\
\textbf{Primal Infaisable} & Impossible & \textcolor{defvert}{Possible} & \textcolor{defvert}{Possible} \\
\textbf{Primal Non Borne} & Impossible & \textcolor{defvert}{Possible} & Impossible \\
\bottomrule
\end{tabular}
\end{center}

\textcolor{importantrouge}{\textbf{Attention :}} Primal infaisable ET dual infaisable est possible simultanement !
\end{important}

\begin{exemple}
\textbf{Exemple d'auto-dualite (Self-Duality) :}

Considerons :
\begin{align*}
\text{minimiser} \quad & z = 4x_1 + 7x_2 + 9x_3\\
\text{sujet a} \quad & x_2 + 2x_3 \geq -4\\
& -x_1 + 6x_3 \geq -7\\
& 2x_1 + 6x_2 \leq 9\\
& x_1, x_2, x_3 \geq 0
\end{align*}

En ecrivant le dual, on obtient un probleme \textcolor{algoviolet}{\textbf{equivalent}} au primal. On parle de probleme \textbf{auto-dual}. Cela se produit lorsque la matrice des contraintes est \textcolor{importantrouge}{\textbf{antisymetrique}} ($A^T = -A$).

\textbf{Consequence :} Primal et dual peuvent etre infaisables en meme temps dans ce cas.
\end{exemple}

%=============================================
\section{Les Ecarts Complementaires}
%=============================================

\needspace{10\baselineskip}
\subsection{Definition des Conditions d'Ecarts Complementaires}

\begin{definition}
\textbf{Conditions d'ecarts complementaires :}

Soient $x^* \in \mathbb{R}^n$ et $y^* \in \mathbb{R}^m$. On dit qu'ils satisfont les \textcolor{defvert}{\textbf{conditions d'ecarts complementaires}} si :

\begin{align}
x_j^* = 0 \quad \text{ou} \quad \sum_{i=1}^{m} a_{ij} y_i^* = c_j \quad \text{(ou les deux)} &\quad \text{pour } j = 1, 2, \ldots, n \tag{1}\\[6pt]
\sum_{j=1}^{n} a_{ij} x_j^* = b_i \quad \text{ou} \quad y_i^* = 0 \quad \text{(ou les deux)} &\quad \text{pour } i = 1, 2, \ldots, m \tag{2}
\end{align}
\end{definition}

\begin{notecours}
\textbf{Interpretation intuitive des deux conditions :}

\begin{itemize}
    \item \textcolor{titrebleu}{\textbf{Condition (1) :}} Pour chaque variable primale, soit elle est nulle ($x_j^* = 0$), soit la variable d'ecart de la contrainte duale correspondante est nulle (la contrainte duale est \textbf{saturee}).
    \item \textcolor{defvert}{\textbf{Condition (2) :}} Pour chaque contrainte primale, soit sa variable d'ecart est nulle (contrainte \textbf{saturee}), soit la variable duale correspondante est nulle ($y_i^* = 0$).
\end{itemize}

En d'autres termes : \textcolor{importantrouge}{\textbf{une variable et son ecart dual ne peuvent etre simultanement strictement positifs.}}
\end{notecours}

\newpage
\subsection{Theoreme des Ecarts Complementaires}

\begin{theoreme}
\textbf{Theoreme 6 (Theoreme des ecarts complementaires) :}

\textit{Soient $x^*$ et $y^*$ des solutions realisables pour les problemes primal et dual, respectivement. Les conditions d'ecarts complementaires sont des conditions \textcolor{defvert}{necessaires et suffisantes} pour l'optimalite simultanee de $x^*$ et $y^*$.}
\end{theoreme}

\begin{methode}
\textbf{Preuve detaillee :}

L'expression suivante (issue de la preuve de la dualite faible) :
$$\sum_{j=1}^n c_j x_j^* \leq \sum_{j=1}^n \left(\sum_{i=1}^m a_{ij} y_i^*\right) x_j^* = \sum_{i=1}^m \left(\sum_{j=1}^n a_{ij} x_j^*\right) y_i^* \leq \sum_{i=1}^m b_i y_i^*$$

est une egalite partout \textcolor{importantrouge}{\textbf{si et seulement si}} :

\begin{align*}
c_j x_j^* &= \left(\sum_{i=1}^m a_{ij} y_i^*\right) x_j^* \quad \text{pour } j = 1, 2, \ldots, n \quad \Rightarrow \textbf{ Condition (1)}\\[6pt]
\left(\sum_{j=1}^n a_{ij} x_j^*\right) y_i^* &= b_i y_i^* \quad \text{pour } i = 1, 2, \ldots, m \quad \Rightarrow \textbf{ Condition (2)}
\end{align*}

Par le theoreme de dualite forte, cette egalite est necessaire et suffisante pour l'optimalite simultanee. $\square$
\end{methode}

\subsection{Forme Plus Applicable}

\begin{theoreme}
\textbf{Theoreme 7 (Forme pratique) :}

\textit{Une solution realisable $x^*$ du probleme primal est optimale si et seulement si il existe des nombres $y_1^*, y_2^*, \ldots, y_m^*$ tels que :}

\textbf{Conditions d'ecarts complementaires :}
\begin{align*}
\sum_{i=1}^{m} a_{ij} y_i^* &= c_j \quad \text{chaque fois que } x_j^* > 0\\
y_i^* &= 0 \quad \text{chaque fois que } \sum_{j=1}^{n} a_{ij} x_j^* < b_i
\end{align*}

\textbf{et $y^*$ est une solution duale :}
\begin{align*}
\sum_{i=1}^{m} a_{ij} y_i^* &\geq c_j \quad \text{pour } j = 1, 2, \ldots, n\\
y_i^* &\geq 0 \quad \text{pour } i = 1, 2, \ldots, m
\end{align*}
\end{theoreme}

\begin{astuce}
\textbf{Usage pratique :}

Ce theoreme permet de \textcolor{defvert}{\textbf{verifier l'optimalite}} d'une solution proposee sans executer la methode du simplexe. La procedure est :

\begin{enumerate}
    \item Verifier que $x^*$ est une solution primale realisable
    \item Utiliser les conditions (1) et (2) pour deduire les valeurs de $y^*$
    \item Verifier que $y^*$ est une solution duale realisable
    \item Si les trois conditions sont remplies : \textcolor{defvert}{\textbf{$x^*$ et $y^*$ sont optimales !}}
\end{enumerate}
\end{astuce}

\newpage
\subsection{Theoreme d'Unicite}

\begin{theoreme}
\textbf{Theoreme 8 (Unicite de la solution duale associee) :}

\textit{Si $x^*$ est une solution de base primale \textcolor{defvert}{non degeneree}, alors le systeme d'equations du Theoreme 7 a une \textcolor{defvert}{solution unique}.}
\end{theoreme}

\begin{important}
\textbf{Implication pratique :}

La strategie de verification par les ecarts complementaires est \textcolor{importantrouge}{\textbf{applicable uniquement}} si le systeme du Theoreme 7 a une solution unique, ce qui est garanti lorsque $x^*$ est non degeneree.

En cas de degenerescence, plusieurs solutions duales peuvent correspondre a la meme solution primale optimale.
\end{important}

%=============================================
\section{Exemples Detailles : Verification par les Ecarts Complementaires}
%=============================================

\needspace{10\baselineskip}
\subsection{Probleme d'Application}

\begin{exemple}
\textbf{Probleme \#4 -- Verification de solutions candidates :}

\begin{align*}
\text{maximiser} \quad & z = -x_1 - 2x_2\\
\text{sujet a} \quad & -3x_1 + x_2 \leq -1\\
& x_1 - x_2 \leq 1\\
& -2x_1 + 7x_2 \leq 6\\
& 9x_1 - 4x_2 \leq 6\\
& -5x_1 + 2x_2 \leq -3\\
& 7x_1 - 3x_2 \leq 6\\
& x_1, x_2 \geq 0
\end{align*}

On pretend que la solution optimale est $x_1^* = \frac{3}{5}$, $x_2^* = 0$.
\end{exemple}

\begin{methode}
\textbf{Construction du dual de ce probleme :}

\begin{align*}
\text{minimiser} \quad & w = -y_1 + y_2 + 6y_3 + 6y_4 - 3y_5 + 6y_6\\
\text{sujet a} \quad & -3y_1 + y_2 - 2y_3 + 9y_4 - 5y_5 + 7y_6 \geq -1\\
& y_1 - y_2 + 7y_3 - 4y_4 + 2y_5 - 3y_6 \geq -2\\
& y_1, y_2, y_3, y_4, y_5, y_6 \geq 0
\end{align*}
\end{methode}

\newpage
\subsection{Premiere Affirmation : Verification par les Ecarts Complementaires}

\begin{exemple}
\textbf{Affirmation 1 :}

$x_1^* = \frac{3}{5}$, $x_2^* = 0$ et $y_1^* = 0$, $y_2^* = 0$, $y_3^* = 0$, $y_4^* = 0$, $y_5^* = \frac{1}{5}$, $y_6^* = 0$ sont des solutions optimales pour le primal et le dual.

\textbf{Tableau de verification :}

\begin{center}
\begin{tabular}{l|c||l|c}
\toprule
\textbf{Variable primale} & \textbf{Valeur} & \textbf{Variable duale} & \textbf{Valeur} \\
\midrule
$x_3$ (ecart contr.1) & $\frac{4}{5}$ & $y_1$ & $0$ \\
$x_4$ (ecart contr.2) & $\frac{2}{5}$ & $y_2$ & $0$ \\
$x_5$ (ecart contr.3) & $\frac{36}{5}$ & $y_3$ & $0$ \\
$x_6$ (ecart contr.4) & $\frac{3}{5}$ & $y_4$ & $0$ \\
$x_7$ (ecart contr.5) & $0$ & $y_5$ & $\frac{1}{5}$ \\
$x_8$ (ecart contr.6) & $\frac{9}{5}$ & $y_6$ & $0$ \\
\midrule
$x_1$ & $\frac{3}{5}$ & $y_7$ (ecart dual c.1) & $0$ \\
$x_2$ & $0$ & $y_8$ (ecart dual c.2) & $\frac{12}{5}$ \\
\midrule
$z$ (primal) & $-\frac{3}{5}$ & $w$ (dual) & $-\frac{3}{5}$ \\
\bottomrule
\end{tabular}
\end{center}

\textbf{Conditions d'ecarts complementaires verifiees :}
\begin{itemize}
    \item Pour chaque contrainte primale non saturee ($x_3, x_4, x_5, x_6, x_8 > 0$), la variable duale correspondante est bien 0 \checkmark
    \item $x_7 = 0$ (contrainte 5 saturee) et $y_5 = \frac{1}{5} > 0$ \checkmark
    \item $x_1 = \frac{3}{5} > 0$ et $y_7 = 0$ (contrainte duale 1 saturee) \checkmark
    \item $x_2 = 0$ et $y_8 > 0$ (la contrainte duale 2 n'est pas saturee) \checkmark
    \item $z^* = w^* = -\frac{3}{5}$ \checkmark
\end{itemize}

\textcolor{defvert}{\textbf{Conclusion : Cette affirmation est VRAIE ! Les deux solutions sont optimales.}}
\end{exemple}
\vspace{0.5cm}
\subsubsection*{Explications Détaillées : De la solution primale à la vérification}

Cette section détaille le raisonnement mathématique caché derrière l'affirmation 1 et la construction du tableau de vérification.

\paragraph{1. Comment a-t-on "deviné" la solution duale candidate ?}
Le théorème des écarts complémentaires nous permet de déduire la solution optimale du dual uniquement à partir de celle du primal ($x_1^* = \frac{3}{5}, x_2^* = 0$), sans avoir à résoudre le dual de zéro. Ce théorème repose sur deux règles fondamentales :

\vspace{0.3cm}
\noindent\textbf{Règle A : Contraintes primales non saturées $\implies$ Variables duales nulles} \\
Si l'on remplace les valeurs de $x^*$ dans les contraintes du primal, on observe quelles contraintes ont un "reste" (non saturées) et lesquelles atteignent exactement leur limite (saturées) :
\begin{itemize}
    \item \textbf{Contrainte 1 :} $-3(\frac{3}{5}) + 0 = -\frac{9}{5}$. Comme $-\frac{9}{5} < -1$, la contrainte n'est pas saturée $\implies y_1 = 0$.
    \item Par le même calcul, les contraintes \textbf{2, 3, 4 et 6} sont strictement inférieures à leur second membre $\implies y_2 = 0, y_3 = 0, y_4 = 0, y_6 = 0$.
    \item \textbf{Contrainte 5 :} $-5(\frac{3}{5}) + 2(0) = -3$. La limite est exactement atteinte ($-3 = -3$). C'est la seule contrainte saturée, donc seule la variable $y_5$ peut être strictement positive.
\end{itemize}

\vspace{0.3cm}
\noindent\textbf{Règle B : Variables primales strictement positives $\implies$ Contraintes duales saturées} \\
Puisque $x_1^* = \frac{3}{5}$, ce qui est strictement supérieur à $0$, la \textbf{première contrainte du dual} (qui correspond à $x_1$) doit obligatoirement être saturée. Le symbole $\geq$ devient une égalité stricte $=$ :
$$ -3y_1 + y_2 - 2y_3 + 9y_4 - 5y_5 + 7y_6 = -1 $$

Puisque nous savons grâce à la Règle A que toutes les variables $y$ (sauf $y_5$) valent $0$, l'équation se simplifie radicalement :
$$ -5y_5 = -1 \implies y_5 = \frac{1}{5} $$
C'est ainsi que l'on obtient la solution duale candidate complète : $y_5^* = \frac{1}{5}$ et tous les autres $y^* = 0$.

\paragraph{2. Comment calculer les valeurs du tableau de vérification ?}
Le tableau met en évidence les "variables d'écart", c'est-à-dire la différence entre les limites imposées par le problème et ce qui est réellement consommé par nos solutions candidates.

\vspace{0.3cm}
\noindent\textbf{Pour le bloc Primal (écarts $x_3$ à $x_8$) :} \\
Les contraintes primales sont de type $\leq$. L'écart se calcule donc par : $\text{Limite (droite)} - \text{Valeur calculée (gauche)}$.
\begin{itemize}
    \item \textit{Exemple pour l'écart de la contrainte 1 ($x_3$) :}
    $$ -1 - \left( -3\left(\frac{3}{5}\right) + 0 \right) = -1 - \left(-\frac{9}{5}\right) = -1 + \frac{9}{5} = \frac{4}{5} $$
    \item \textit{Exemple pour l'écart de la contrainte 5 ($x_7$) :}
    $$ -3 - \left( -5\left(\frac{3}{5}\right) + 0 \right) = -3 - (-3) = 0 $$
    (L'écart est bien nul car la contrainte est saturée).
\end{itemize}

\vspace{0.3cm}
\noindent\textbf{Pour le bloc Dual (écarts $y_7$ à $y_8$) :} \\
Les contraintes duales sont de type $\geq$. L'écart se calcule dans le sens inverse : $\text{Valeur calculée (gauche)} - \text{Limite (droite)}$.
\begin{itemize}
    \item \textit{Exemple pour l'écart dual de la contrainte 2 ($y_8$) :} En utilisant $y_5 = \frac{1}{5}$ (et les autres à $0$), on évalue la partie gauche de la 2ème contrainte duale : $2\left(\frac{1}{5}\right) = \frac{2}{5}$. \\
    L'écart est donc :
    $$ \frac{2}{5} - (-2) = \frac{2}{5} + \frac{10}{5} = \frac{12}{5} $$
\end{itemize}

\vspace{0.3cm}
\noindent La conclusion du tableau confirme la théorie : à chaque fois qu'une variable est strictement positive (comme $x_1$ ou $y_5$), la variable d'écart qui lui fait face de l'autre côté du miroir vaut exactement $0$.
\newpage
\subsection{Deuxieme Affirmation : Solution Duale Infaisable}

\begin{exemple}
\textbf{Affirmation 2 :}

$x_1^* = \frac{3}{5}$, $x_2^* = 0$ et $y_1^* = \frac{1}{5}$, $y_2^* = 0$, $y_3^* = 0$, $y_4^* = 0$, $y_5^* = \frac{2}{15}$, $y_6^* = 0$ sont des solutions optimales.

\textbf{Verification :}

En calculant la valeur de la variable d'ecart $y_7$ de la premiere contrainte duale :
$$y_7 = -3y_1 + y_2 - 2y_3 + 9y_4 - 5y_5 + 7y_6 - (-1) = -\frac{3}{5} \cdot \frac{1}{5} - \frac{5 \cdot 2}{15} + 1 = -\frac{4}{15} < 0$$

\textcolor{importantrouge}{\textbf{La premiere contrainte duale n'est pas satisfaite !}} $y_7 = -\frac{4}{15} < 0$.

\textcolor{importantrouge}{\textbf{Conclusion : Cette affirmation est FAUSSE. La solution duale proposee n'est pas realisable.}}
\end{exemple}

\subsection{Troisieme Affirmation : Solution Primale Infaisable}

\begin{exemple}
\textbf{Affirmation 3 :}

$x_1^* = \frac{1}{5}$, $x_2^* = \frac{1}{5}$ et $y_1^* = 0$, $y_2^* = 0$, $y_3^* = 0$, $y_4^* = 0$, $y_5^* = \frac{1}{5}$, $y_6^* = 0$ sont des solutions optimales.

\textbf{Verification :}

Calcul de la variable d'ecart de la premiere contrainte primale :
$$x_3 = -1 - (-3 \cdot \frac{1}{5} + \frac{1}{5}) = -1 + \frac{2}{5} = -\frac{3}{5} < 0$$

\textcolor{importantrouge}{\textbf{La premiere contrainte primale n'est pas satisfaite !}}

Calcul de la variable d'ecart de la cinquieme contrainte primale :
$$x_7 = -3 - (-5 \cdot \frac{1}{5} + 2 \cdot \frac{1}{5}) = -3 + 1 - \frac{2}{5} = -\frac{12}{5} < 0$$

\textcolor{importantrouge}{\textbf{Conclusion : Cette affirmation est FAUSSE. La solution primale proposee n'est pas realisable.}}
\end{exemple}

%=============================================
\section{Resume et Points Cles}
%=============================================

\begin{important}
\textbf{Recapitulatif de la theorie de la dualite :}

\begin{enumerate}
    \item \textcolor{titrebleu}{\textbf{Construction du dual :}}
    \begin{itemize}
        \item Utiliser le tableau de correspondance primal/dual
        \item Maximisation $\leftrightarrow$ minimisation ; contrainte $\leq$ $\leftrightarrow$ variable $\geq 0$
        \item Le dual du dual est le primal
    \end{itemize}
    
    \item \textcolor{defvert}{\textbf{Dualite faible :}}
    \begin{itemize}
        \item Toute solution primale $\leq$ toute solution duale
        \item Si les deux valeurs sont egales, les deux solutions sont optimales
    \end{itemize}
    
    \item \textcolor{algoviolet}{\textbf{Dualite forte :}}
    \begin{itemize}
        \item Si le primal est optimal, le dual est optimal, et les valeurs optimales sont egales
        \item La solution duale optimale se lit dans le dictionnaire primal optimal
    \end{itemize}
    
    \item \textcolor{importantrouge}{\textbf{Ecarts complementaires :}}
    \begin{itemize}
        \item Outil de verification d'optimalite sans relancer le simplexe
        \item Conditions necessaires ET suffisantes pour l'optimalite simultanee
    \end{itemize}
\end{enumerate}
\end{important}

\begin{astuce}
\textbf{Retenir le tableau des combinaisons :}

\begin{center}
\begin{tabular}{l|ccc}
\toprule
 & Dual Optimal & Dual Infaisable & Dual Non Borne \\
\midrule
\textbf{Primal Optimal} & \textcolor{defvert}{\checkmark} & \ding{55} & \ding{55} \\
\textbf{Primal Infaisable} & \ding{55} & \textcolor{defvert}{\checkmark} & \textcolor{defvert}{\checkmark} \\
\textbf{Primal Non Borne} & \ding{55} & \textcolor{defvert}{\checkmark} & \ding{55} \\
\bottomrule
\end{tabular}
\end{center}

\textcolor{importantrouge}{\textbf{Rappel :}} Les deux problemes peuvent etre infaisables simultanement !
\end{astuce}

\newpage
%=============================================
\section{Exercices}
%=============================================

\begin{exemple}
\textbf{Exercice 1 -- Ecriture du dual}

Ecrire le dual du probleme primal suivant :
\begin{align*}
\text{maximiser} \quad & z = 2x_1 + 3x_2 - x_3\\
\text{sujet a} \quad & x_1 + x_2 + x_3 \leq 10\\
& 2x_1 - x_2 + 3x_3 \leq 15\\
& x_1, x_2, x_3 \geq 0
\end{align*}

\begin{enumerate}
    \item Ecrire le dual
    \item Verifier que le dual du dual est bien le primal
    \item Si on vous donne la solution primale $x_1^* = 5, x_2^* = 5, x_3^* = 0$, utiliser les ecarts complementaires pour trouver la solution duale optimale
\end{enumerate}
\end{exemple}

\begin{exemple}
\textbf{Exercice 2 -- Correspondance primal/dual avancee}

Ecrire le dual des problemes suivants :

\textbf{a)}
\begin{align*}
\text{minimiser} \quad & z = 3x_1 + 5x_2\\
\text{sujet a} \quad & x_1 + 2x_2 \geq 4\\
& 2x_1 + x_2 \geq 3\\
& x_1, x_2 \geq 0
\end{align*}

\textbf{b)}
\begin{align*}
\text{maximiser} \quad & z = x_1 + 2x_2\\
\text{sujet a} \quad & x_1 + x_2 = 5\\
& x_1 - x_2 \leq 3\\
& x_1 \geq 0, \quad x_2 \text{ libre}
\end{align*}
\end{exemple}

\begin{exemple}
\textbf{Exercice 3 -- Verification par les ecarts complementaires}

Considerons le probleme :
\begin{align*}
\text{maximiser} \quad & z = 5x_1 + 4x_2\\
\text{sujet a} \quad & 6x_1 + 4x_2 \leq 24\\
& x_1 + 2x_2 \leq 6\\
& x_1, x_2 \geq 0
\end{align*}

Quelqu'un affirme que $x_1^* = 3, x_2^* = \frac{3}{2}$ est la solution optimale.

\begin{enumerate}
    \item Ecrire le dual de ce probleme
    \item En utilisant les conditions d'ecarts complementaires, trouver la solution duale correspondante
    \item Verifier que cette solution duale est realisable
    \item Conclure sur l'optimalite de $x^*$
\end{enumerate}
\end{exemple}

\begin{exemple}
\textbf{Exercice 4 -- Analyse de combinaisons primal/dual}

Pour chacun des cas suivants, indiquer si la situation est \textbf{possible ou impossible}, et justifier :

\begin{enumerate}
    \item Le primal a une solution optimale de valeur 10 et le dual a une solution optimale de valeur 8.
    \item Le primal est non borne et le dual a une solution optimale de valeur 5.
    \item Le primal est infaisable et le dual est infaisable.
    \item Le primal a une solution realisable de valeur 7 et le dual a une solution realisable de valeur 9.
    \item Le primal est non borne et le dual est non borne.
\end{enumerate}
\end{exemple}

\begin{exemple}
\textbf{Exercice 5 -- Lecture de la solution duale dans le dictionnaire}

Apres resolution par le simplexe, on obtient le dictionnaire optimal suivant pour un probleme a 3 contraintes et 2 variables de decision ($n=2$, $m=3$) :
\begin{align*}
x_1 &= 4 - 2x_3 + x_5\\
x_2 &= 1 + x_3 - x_4\\
x_3 &= 2 - x_4 + 3x_5 - x_6\\
z &= 15 - 3x_4 - 2x_5 - 4x_6
\end{align*}

ou $x_3, x_4, x_5$ sont les variables d'ecart primales et $x_6$ est la troisieme variable d'ecart.

\begin{enumerate}
    \item Donner la solution primale optimale
    \item Lire la solution duale optimale $y_1^*, y_2^*, y_3^*$ dans ce dictionnaire
    \item Verifier les conditions d'ecarts complementaires
\end{enumerate}
\end{exemple}

\end{document}
